\documentclass[../../../include/open-logic-section]{subfiles}

\begin{document}

\olfileid{sth}{spine}{valpha}

\olsection{تعريف المراحل بوصفها مجموعات~$V_\alpha$}

بعد تعريف الترتيبات الجيدة وأعداد فون نويمان الترتيبية في
\olref[sth][ordinals][]{chap}، نستعين بها هنا على وصف هرم المجموعات
\emph{نفسه}. ونستحضر لذلك العدد الترتيبي الخلفي والعدد الحدي، وقد
تقدم تعريفهما في~\olref[sth][ordinals][opps]{sec}؛ فعليهما يقوم
التعريف الآتي.

\begin{defn}\ollabel{defValphas}
\begin{align*}
V_\emptyset &\defis \emptyset\\
V_{\ordsucc{\alpha}} &\defis \Pow{V_\alpha} & &
\text{لكل عدد ترتيبي }\alpha\\
V_{\alpha} &\defis \bigcup_{\gamma < \alpha} V_\gamma & &
\text{عندما يكون }\alpha\text{ عددًا ترتيبيًا حديًا}
\end{align*}
\end{defn}
\noindent
هذا تعريف بـ\emph{العودية المتناهية التجاوز} على الأعداد الترتيبية.
ولنقابله بالتعريف العودي للدوال على الأعداد الطبيعية.
\footnote{قارن تعريف الجمع والضرب والأس في~\olref[sfr][infinite][dedekind]{sec}.}
فهناك وهنا نعين حالة الأساس والحالات الخلفية؛ غير أن الأعداد
الترتيبية تستدعي، فوق ذلك، بيان ما يقع في الحالات \emph{الحدية}.

ويمثل تعريف مجموعات~$V_\alpha$ خطوة أساسية. فقد مهدنا لهرم
المجموعات بتصور غير صوري تتكون فيه المجموعات على \emph{مراحل}؛
ومجموعات~$V_\alpha$ هي تلك المراحل في الحقيقة. وإنما تجدر العناية
بأن وصف المراحل هنا \emph{داخلي}: لا نقترح \emph{نموذجًا} محتملًا
للنظرية، بل نعرّف المراحل \emph{في داخل} نظرية المجموعات نفسها.

\end{document}
